\chapter{Other Estimators and Aliases}

\section{Commands not covered in detail}

This chapter collects additional commands and aliases that are
implemented in \hayashi{} but not treated in a dedicated chapter.

\section{Additional discrete and panel models}

\begin{itemize}
\item \texttt{cmnlogit}: conditional multinomial logit
\item \texttt{clogit}: conditional logit
\item \texttt{cpoisson}: conditional Poisson / PPML
\item \texttt{panel\_tobit}: panel Tobit with random effects
\item \texttt{panel\_heckman}: panel Heckman selection
\item \texttt{panel\_qreg}: panel quantile regression
\item \texttt{sysgmm}: system GMM
\item \texttt{pthresh}: panel threshold model
\end{itemize}

\section{Additional time series}

\begin{itemize}
\item \texttt{varma} / \texttt{varmax}: VARMA/VARMAX
\item \texttt{svar\_irf} / \texttt{sirf}: structural VAR impulse responses
\item \texttt{svar\_fevd} / \texttt{sfevd}: structural forecast-error variance
      decomposition
\item \texttt{qvar}: quantile VAR
\item \texttt{tvar}: threshold VAR
\item \texttt{mstl}: multiple-seasonal STL decomposition
\item \texttt{modwt}: maximal overlap discrete wavelet transform
\item \texttt{decompose}: classical time-series decomposition
\item \texttt{dfm}: dynamic factor model
\item \texttt{midas}: mixed-data-sampling regression
\item \texttt{hawkes}: self-exciting point process
\end{itemize}

\section{Additional spatial variants}

\begin{itemize}
\item \texttt{spatial\_sar}, \texttt{spatial\_sem}, \texttt{spatial\_durbin},
      \texttt{spatial\_panel\_sar}, \texttt{spatial\_panel\_sem},
      \texttt{spatial\_durbin\_error}: the spatial variants are also
      available under the aliases \texttt{sdm} and \texttt{sdem}.
\item \texttt{fuzzy\_rd}: fuzzy regression discontinuity
\end{itemize}

\section{Additional tests and diagnostics}

\begin{itemize}
\item \texttt{breusch\_godfrey} / \texttt{bg} / \texttt{bgodfrey}:
      Breusch-Godfrey serial-correlation LM test
\item \texttt{durbinwatson} / \texttt{dw}: Durbin-Watson test
\item \texttt{phillips\_perron} / \texttt{pp}: Phillips-Perron unit-root test
\item \texttt{zivot\_andrews} / \texttt{za}: Zivot-Andrews unit-root test with
      a structural break
\item \texttt{engle\_granger}: Engle-Granger cointegration test
\item \texttt{anderson\_darling} / \texttt{adtest}: Anderson-Darling normality
      test
\item \texttt{harvey\_collier}: Harvey-Collier misspecification test
\item \texttt{lilliefors}: Lilliefors normality test
\item \texttt{sktest}: skewness/kurtosis test
\item \texttt{spearman}, \texttt{wilcoxon\_rs}, \texttt{wilcoxon\_signed\_rank},
      \texttt{kruskal\_wallis}: non-parametric tests
\item \texttt{proptest}, \texttt{proptest2}, \texttt{propci}, \texttt{chi2\_2x2},
      \texttt{multipletests}: proportion and multiple-testing commands
\end{itemize}

\section{Post-estimation tests}

\begin{itemize}
\item \texttt{estat\_overid} / \texttt{sargan}: Sargan/Hansen J over-
      identification test
\item \texttt{estat\_endog} / \texttt{dwh}: Durbin-Wu-Hausman endogeneity test
\item \texttt{estat\_classification}: classification table for logit/probit
\item \texttt{estat\_gof} / \texttt{hltest}: Hosmer-Lemeshow goodness-of-fit
\end{itemize}

\section{Aliases}

Many estimators have aliases for convenience. The most common are:

\begin{itemize}
\item \texttt{reg} for \texttt{ols}, \texttt{frontier} for \texttt{sfa\_production}
\item \texttt{neural\_net} for \texttt{mlp}, \texttt{xgb} for \texttt{xgboost}
\item \texttt{causal\_forest} for \texttt{causalforest}, \texttt{random\_forest} for \texttt{rf}
\item \texttt{gaussian\_process} for \texttt{gp}, \texttt{bayesian\_trees} for \texttt{bart}
\item \texttt{gmm\_clustering} for \texttt{gmm\_clust}, \texttt{pca\_biplot} for \texttt{biplot}
\item \texttt{three\_sls}, \texttt{threesl}, \texttt{3sls}, \texttt{reg3}: aliases for the
      same 3SLS estimator
\end{itemize}

Use the quick-reference appendix for the full matrix of commands.
